\section{Conclusion}\label{sec:conclusion}
\begin{CJK*}{UTF8}{gbsn}

\methodrevision{本文围绕结构化知识如何适配当前音频并有效参与零样本分类的问题，提出了样本自适应知识整合框架SAKI。该框架通过样本级关系选择确定参与知识增强的关系，利用AAKV将图谱三元组转换为受约束的声音描述，并采用分层证据融合协调知识证据与原始CLAP预测对最终类别排序的贡献。在五个测试集上，SAKI相较于本文复现的iKnow-audio将Hit@1提高了2.08--12.73个百分点，MRR提高了1.02--6.68个百分点。完整因子消融、关系选择对照和知识语言化实验表明，三个模块在五数据集等权平均指标上具有互补作用，但单项数据集的增益并非完全一致。}

本文结果说明，结构化知识在零样本音频分类中的作用不仅取决于图谱是否包含相关实体，还取决于所选知识关系是否能够为当前样本提供有效的判别信息、关系语义能否被音频–文本模型有效表示，以及知识分数能否以适当方式参与最终决策。SAKI为将音频中心知识图谱转化为样本相关的分类证据提供了一种可行方案。当前结论适用于本文采用的CLAP模型、RotatE知识图谱表示及候选关系池；未来可通过证据可信度评估、关系预筛选以及跨模型和跨知识图谱验证，进一步提高知识增强分类的可靠性与适用范围。

\end{CJK*}
